\documentclass[12pt]{article}
\usepackage[utf8]{inputenc}
\usepackage[T1]{fontenc}
\usepackage{hyperref}
\title{Tasawaq Language Resources for K-12 Curriculum Development}
\author{Compiled by Claude Sonnet 5 (claude-sonnet-5) for Andr\'as Kornai}
\date{}
\begin{document}
\maketitle
\section{Introduction}
Tasawaq (ISO 639-3 twq) is a Northern Songhay language of the In-Gall area of Agadez region, Niger, with strong Tamajaq influence. Speaker estimates differ (about 8,000 in the survey data; Wikipedia gives 21,000). Searches found no school materials in Tasawaq; what exists is mainly linguistic description.
\section{Primary Educational Resources}
No primers, textbooks or school programmes for Tasawaq were found in search results.
\section{Linguistic and Reference Materials}
\begin{itemize}
\item Wikipedia overview of the Tasawaq language: \url{https://en.wikipedia.org/wiki/Tasawaq_language}
\item Lexilogos Songhay dictionary page, which lists a thematic Tasawaq--French dictionary for In-Gall with audio (its contents were not opened): \url{https://www.lexilogos.com/songhai_dictionnaire.htm}
\item Jeffrey Heath's Songhay languages page (University of Michigan): \url{https://websites.umich.edu/~jheath/Songhay%20languages.html}
\item Leiden scholarly publication on Songhay grammar, cited in the search results: \url{https://scholarlypublications.universiteitleiden.nl/access/item%3A2937663/view}
\end{itemize}
Search results also mention an unpublished 1988 Niamey MA thesis by Ousseina Alidou, a linguistic sketch of Tasawaq d'In-Gall. No online copy was found.
\section{Religious and Community Resources}
None found.
\section{Digital and Audio Resources}
The audio mentioned on the Lexilogos page is the only one found, and it was not verified.
\section{Curriculum Development Context}
Very little is available. A curriculum would have to be built from fieldwork and the linguistic descriptions above, with Tamajaq, Hausa and French as contact languages.
\end{document}